\documentclass[12pt]{article}
\usepackage[a4paper,margin=2.2cm]{geometry}
\usepackage{amsmath,amssymb,amsthm,mathtools}
\usepackage{enumitem}
\usepackage{hyperref}
\usepackage{ctex} % 中文支持

\title{第10章：向量丛与纤维丛}
\author{}
\date{}

\newtheorem{definition}{定义}[section]
\newtheorem{theorem}[definition]{定理}
\newtheorem{lemma}[definition]{引理}
\newtheorem{proposition}[definition]{命题}
\newtheorem{corollary}[definition]{推论}
\newtheorem{example}[definition]{例}
\newtheorem{exercise}[definition]{练习}

\begin{document}
\maketitle

\tableofcontents

\section{向量丛：基本定义与例子}

\begin{definition}[向量丛（Vector bundle）]
设 $M$ 为拓扑空间。一个\textbf{秩为 $k$ 的（实）向量丛}是一个拓扑空间 $E$ 连同一个满射连续映射
\[
\pi:E\to M
\]
满足：
\begin{enumerate}[label=(\roman*)]
\item 对每个 $p\in M$，纤维 $E_p=\pi^{-1}(p)$ 赋予 $k$ 维实向量空间结构。
\item 对每个 $p\in M$，存在 $p$ 的邻域 $U\subset M$ 与同胚
\[
\Phi:\pi^{-1}(U)\to U\times \mathbb{R}^k
\]
（称为 $U$ 上的\textbf{局部平凡化}），使得
\[
\pi_1\circ \Phi=\pi,\quad \text{且对每个 }q\in U,\ \Phi|_{E_q}:E_q\to \{q\}\times \mathbb{R}^k\cong \mathbb{R}^k
\]
为向量空间同构。
\end{enumerate}
若 $M,E$ 为（可带边界的）光滑流形，$\pi$ 光滑，且可取 $\Phi$ 为微分同胚，则称为\textbf{光滑向量丛}。
\end{definition}

\begin{definition}[总空间、底空间、投影、线丛、平凡丛]
$E$ 称为\textbf{总空间}，$M$ 称为\textbf{底空间}，$\pi$ 称为\textbf{投影}。
秩为 $1$ 的向量丛称为\textbf{线丛}。
若存在全局平凡化（全体 $M$ 上的一个局部平凡化），则称 $E$ 为\textbf{平凡丛}；在光滑情形若全局平凡化为微分同胚，则称\textbf{光滑平凡}。
\end{definition}

\begin{exercise}[10.1]
设 $E$ 是 $M$ 上光滑向量丛。证明投影 $\pi:E\to M$ 是满射的光滑下沉映射（surjective smooth submersion）。
\end{exercise}

\begin{example}[10.2（积丛 / 直积丛）]
任意拓扑空间 $M$ 上，$E=M\times \mathbb{R}^k$ 配以投影 $\pi=\pi_1$ 为秩 $k$ 向量丛，称为\textbf{积丛（product bundle）}，是平凡丛。
\end{example}

\begin{example}[10.3（莫比乌斯丛）]
在 $\mathbb{R}^2$ 上定义等价关系 $(x,y)\sim(x+n,(-1)^n y)$（$n\in\mathbb{Z}$），令 $E=\mathbb{R}^2/\sim$，商映射 $q:\mathbb{R}^2\to E$。
令覆盖映射 $\varepsilon:\mathbb{R}\to S^1$，$\varepsilon(x)=e^{2\pi i x}$，则存在连续映射 $\pi:E\to S^1$ 使得 $\pi\circ q=\varepsilon\circ \pi_1$。
在适当光滑结构下，$\pi:E\to S^1$ 是 $S^1$ 上的光滑实线丛，称为\textbf{莫比乌斯丛}。
\end{example}

\subsection{切丛作为向量丛}

\begin{proposition}[10.4（切丛是向量丛）]
设 $M$ 是（可带边界的）光滑 $n$ 维流形，$TM$ 为其切丛。赋予 $TM$ 标准投影、纤维上的自然向量空间结构，以及自然坐标构造的拓扑与光滑结构，则 $TM$ 是 $M$ 上秩 $n$ 的光滑向量丛。
\end{proposition}

\begin{lemma}[10.5（过渡函数的形式）]
设 $\pi:E\to M$ 为秩 $k$ 光滑向量丛，$\Phi:\pi^{-1}(U)\to U\times\mathbb{R}^k$、$\Psi:\pi^{-1}(V)\to V\times\mathbb{R}^k$ 为光滑局部平凡化，且 $U\cap V\neq\varnothing$。
则存在光滑映射 $\tau:U\cap V\to GL(k,\mathbb{R})$ 使得
\[
\Phi\circ\Psi^{-1}(p,v)=(p,\tau(p)v),\qquad (p,v)\in (U\cap V)\times\mathbb{R}^k.
\]
$\tau$ 称为（这两平凡化之间的）\textbf{过渡函数}。
\end{lemma}

\begin{lemma}[10.6（向量丛图册引理 / Chart Lemma）]
设 $M$ 为（可带边界的）光滑流形，并对每个 $p\in M$ 给定固定维数 $k$ 的实向量空间 $E_p$。令
\[
E=\bigsqcup_{p\in M}E_p,\qquad \pi:E\to M
\]
为自然投影。设给定
\begin{enumerate}[label=(\roman*)]
\item 开覆盖 $\{U_\alpha\}$；
\item 双射 $\Phi_\alpha:\pi^{-1}(U_\alpha)\to U_\alpha\times\mathbb{R}^k$，且在每个纤维上线性同构；
\item 对任意 $U_\alpha\cap U_\beta\neq\varnothing$，存在光滑 $\tau_{\alpha\beta}:U_\alpha\cap U_\beta\to GL(k,\mathbb{R})$ 使
\[
\Phi_\alpha\circ\Phi_\beta^{-1}(p,v)=(p,\tau_{\alpha\beta}(p)v).
\]
\end{enumerate}
则 $E$ 上存在唯一拓扑与光滑结构，使 $\pi:E\to M$ 成为秩 $k$ 光滑向量丛投影且 $\{(U_\alpha,\Phi_\alpha)\}$ 为光滑局部平凡化。
\end{lemma}

\subsection{由已知向量丛构造新向量丛}

\begin{example}[10.7（Whitney 和 / 直和丛）]
若 $E'\to M$、$E''\to M$ 为秩分别 $k',k''$ 的光滑向量丛，则其\textbf{Whitney 和}定义为
\[
E'\oplus E''=\bigsqcup_{p\in M}(E'_p\oplus E''_p)\to M,
\]
它是秩 $k'+k''$ 的光滑向量丛；过渡函数为块对角矩阵 $\tau'(p)\oplus\tau''(p)$。
\end{example}

\begin{example}[10.8（限制丛 / Restriction）]
设 $E\to M$ 为向量丛，$S\subseteq M$。定义限制
\[
E|_S=\bigcup_{p\in S}E_p,\qquad E|_S\to S
\]
为投影限制。若 $E$ 为光滑丛且 $S$ 为浸入/嵌入子流形，则 $E|_S$ 为光滑向量丛。特别地 $TM|_S$ 称为\textbf{环境切丛（ambient tangent bundle）}。
\end{example}

\section{截面（Sections）与截面空间}

\begin{definition}[（局部/整体）截面与粗糙截面]
设 $\pi:E\to M$ 为向量丛。
\begin{itemize}
\item \textbf{整体截面}：连续映射 $\sigma:M\to E$ 满足 $\pi\circ\sigma=\mathrm{Id}_M$。
\item \textbf{局部截面}：在开集 $U\subseteq M$ 上连续映射 $\sigma:U\to E$ 满足 $\pi\circ\sigma=\mathrm{Id}_U$。
\item \textbf{光滑截面}：若 $E,M$ 光滑，且 $\sigma$ 为光滑映射。
\item \textbf{粗糙截面（rough section）}：同样满足 $\pi\circ\sigma=\mathrm{Id}$，但不要求连续。
\end{itemize}
\end{definition}

\begin{definition}[零截面与支集]
\textbf{零截面} $\zeta:M\to E$ 定义为 $\zeta(p)=0\in E_p$。
截面 $\sigma$ 的\textbf{支集}为集合 $\{p:\sigma(p)\neq 0\}$ 的闭包。
\end{definition}

\begin{exercise}[10.9]
证明：任意向量丛的零截面连续；任意光滑向量丛的零截面光滑。
\end{exercise}

\begin{example}[10.10（截面的例子）]
设 $M$ 为光滑流形：
\begin{enumerate}[label=(\alph*)]
\item $TM$ 的截面就是向量场。
\item 若 $S\subseteq M$ 为浸入子流形，则 $TM|_S\to S$ 的截面称为\textbf{沿 $S$ 的向量场}（与 $TS$ 的向量场不同）。
\item 平凡丛 $M\times\mathbb{R}^k$ 的截面与连续映射 $F:M\to\mathbb{R}^k$ 一一对应：$\tilde F(x)=(x,F(x))$。
\item 特别地，$C^\infty(M)$ 与平凡线丛 $M\times\mathbb{R}$ 的光滑截面空间自然同一。
\end{enumerate}
\end{example}

\begin{definition}[光滑截面空间与模结构]
光滑向量丛 $E\to M$ 的光滑整体截面集合记为 $\Gamma(E)$。
在逐点加法与数乘下它是向量空间，并且对 $f\in C^\infty(M)$、$\sigma\in\Gamma(E)$，定义
\[
(f\sigma)(p)=f(p)\sigma(p),
\]
则 $\Gamma(E)$ 是 $C^\infty(M)$ 上的模。
\end{definition}

\begin{exercise}[10.11]
设 $E\to M$ 光滑向量丛：
\begin{enumerate}[label=(\alph*)]
\item 若 $\sigma,\tau\in\Gamma(E)$，$f,g\in C^\infty(M)$，证明 $f\sigma+g\tau\in\Gamma(E)$；
\item 证明 $\Gamma(E)$ 是 $C^\infty(M)$-模。
\end{enumerate}
\end{exercise}

\begin{lemma}[10.12（向量丛的延拓引理 / Extension Lemma）]
设 $\pi:E\to M$ 为光滑向量丛，$A\subseteq M$ 为闭集，$\sigma:A\to E$ 为限制丛 $E|_A$ 的截面，且在 $A$ 的每点附近可延拓为光滑局部截面。
则对任意包含 $A$ 的开集 $U\subseteq M$，存在整体光滑截面 $\tilde\sigma\in\Gamma(E)$ 使
\[
\tilde\sigma|_A=\sigma,\qquad \mathrm{supp}(\tilde\sigma)\subseteq U.
\]
\end{lemma}

\begin{exercise}[10.13]
证明引理 10.12。
\end{exercise}

\begin{exercise}[10.14]
设 $\pi:E\to M$ 为光滑向量丛。证明：$E$ 的每个元素都在某个光滑整体截面的像中。
\end{exercise}

\section{标架（Frames）与平凡化}

\begin{definition}[线性无关、张成、局部/整体标架]
设 $E\to M$ 秩为 $k$，$U\subseteq M$ 为开集。
\begin{itemize}
\item $k$ 个局部截面 $(\sigma_1,\dots,\sigma_k)$ 在 $U$ 上\textbf{线性无关}：对每个 $p\in U$，$\{\sigma_i(p)\}$ 在线性代数意义下无关。
\item \textbf{张成}：对每个 $p$，$\{\sigma_i(p)\}$ 张成 $E_p$。
\item \textbf{局部标架（local frame）}：在 $U$ 上线性无关且张成的有序 $k$-元组；若 $U=M$ 称\textbf{整体标架（global frame）}。
\item 若各 $\sigma_i$ 光滑，则称\textbf{光滑标架}。
\end{itemize}
\end{definition}

\begin{proposition}[10.15（补全局部标架）]
设 $E\to M$ 为秩 $k$ 光滑向量丛：
\begin{enumerate}[label=(\alph*)]
\item 若 $(\sigma_1,\dots,\sigma_m)$ 为 $U$ 上线性无关的光滑局部截面（$1\le m<k$），则对每个 $p\in U$，存在邻域 $V$ 及光滑截面 $\sigma_{m+1},\dots,\sigma_k$ 使得在 $U\cap V$ 上构成光滑局部标架。
\item 若 $(v_1,\dots,v_m)$ 为 $E_p$ 中线性无关的 $m$-元组（$1\le m\le k$），则存在 $p$ 邻域上的光滑局部标架 $(\sigma_i)$ 满足 $\sigma_i(p)=v_i$（$i\le m$）。
\item 若 $A\subseteq M$ 闭，且 $E|_A$ 上给定线性无关的 $k$ 个截面 $(\tau_1,\dots,\tau_k)$（以引理 10.12 的意义局部可光滑延拓），则在 $A$ 的某邻域存在光滑局部标架 $(\sigma_i)$ 满足 $\sigma_i|_A=\tau_i$。
\end{enumerate}
\end{proposition}

\begin{exercise}[10.16]
证明命题 10.15。
\end{exercise}

\begin{example}[10.17（积丛的整体标架）]
对平凡丛 $E=M\times\mathbb{R}^k$，$\mathbb{R}^k$ 的标准基 $(e_i)$ 给出整体标架 $\tilde e_i(p)=(p,e_i)$（光滑情形为光滑整体标架）。
\end{example}

\begin{example}[10.18（由平凡化诱导标架）]
若 $\Phi:\pi^{-1}(U)\to U\times\mathbb{R}^k$ 为局部平凡化，则定义
\[
\sigma_i(p)=\Phi^{-1}(p,e_i)
\]
得到 $U$ 上的局部标架，称与 $\Phi$ \textbf{相联系}。
\end{example}

\begin{proposition}[10.19（由标架诱导平凡化）]
光滑向量丛的每个光滑局部标架都与某个光滑局部平凡化相联系。更具体地，若 $(\sigma_i)$ 为 $U$ 上局部标架，则
\[
\Psi:U\times\mathbb{R}^k\to \pi^{-1}(U),\quad \Psi(p,(v^1,\dots,v^k))=v^i\sigma_i(p)
\]
为微分同胚，其逆给出相应局部平凡化。
\end{proposition}

\begin{corollary}[10.20（平凡化与整体标架）]
光滑向量丛\textbf{光滑平凡}当且仅当它承认一个\textbf{光滑整体标架}。
特别地，$TM$ 平凡当且仅当 $M$ 可平行化（parallelizable）。
\end{corollary}

\begin{corollary}[10.21（由局部标架构造总空间坐标图）]
设 $E\to M$ 为秩 $k$ 光滑向量丛，$(V,\varphi)$ 为 $M$ 的坐标图，且 $E$ 在 $V$ 上有光滑局部标架 $(\sigma_i)$。
定义
\[
\tilde\varphi:\pi^{-1}(V)\to \varphi(V)\times\mathbb{R}^k,\quad 
\tilde\varphi\bigl(v^i\sigma_i(p)\bigr)=\bigl(x^1(p),\dots,x^n(p),v^1,\dots,v^k\bigr),
\]
则 $(\pi^{-1}(V),\tilde\varphi)$ 是 $E$ 的光滑坐标图。
\end{corollary}

\begin{proposition}[10.22（截面光滑性的分量判别）]
设 $E\to M$ 光滑向量丛，$\tau$ 为 $U\subseteq M$ 上粗糙截面，$(\sigma_i)$ 为 $U$ 上光滑局部标架。
将 $\tau$ 唯一表示为
\[
\tau(p)=\tau^i(p)\sigma_i(p),
\]
则 $\tau$ 在 $U$ 上光滑当且仅当每个分量函数 $\tau^i:U\to\mathbb{R}$ 光滑。
\end{proposition}

\begin{exercise}[10.23]
证明：粗糙截面连续当且仅当它在每个局部标架中的分量函数都连续。
\end{exercise}

\begin{proposition}[10.24（切丛光滑结构的唯一性）]
设 $M$ 为光滑 $n$ 维流形。命题 3.18 构造的 $TM$ 上拓扑与光滑结构，是唯一使得
\[
\pi:TM\to M
\]
成为光滑向量丛且所有坐标向量场均为光滑局部截面的结构。
\end{proposition}

\section{丛同态（Bundle Homomorphisms）}

\begin{definition}[丛同态与覆盖映射]
设 $\pi:E\to M$、$\pi':E'\to M'$ 为向量丛。若存在映射 $f:M\to M'$ 使得
\[
\pi'\circ F=f\circ \pi,
\]
且对每个 $p\in M$，限制 $F|_{E_p}:E_p\to E'_{f(p)}$ 为线性映射，则称 $F:E\to E'$ 为\textbf{丛同态}，并称 $F$ \textbf{覆盖} $f$。
\end{definition}

\begin{proposition}[10.25]
设 $F:E\to E'$ 为覆盖 $f:M\to M'$ 的丛同态，则 $f$ 连续且由 $F$ 唯一确定；事实上
\[
f=\pi'\circ F\circ \zeta
\]
其中 $\zeta$ 为 $E$ 的零截面。若 $E,E',F$ 光滑，则 $f$ 光滑。
\end{proposition}

\begin{definition}[丛同构、在 $M$ 上的丛同态]
双射丛同态 $F$ 若其逆亦为丛同态，则称为\textbf{丛同构}；若 $F$ 为微分同胚则称\textbf{光滑丛同构}。
当 $E,E'$ 同底空间 $M$ 时，覆盖恒等映射的丛同态称为\textbf{在 $M$ 上的丛同态}，即满足 $\pi'\circ F=\pi$。
\end{definition}

\begin{proposition}[10.26]
设 $E,E'$ 为同一底流形 $M$ 上的光滑向量丛，$F:E\to E'$ 为在 $M$ 上的双射光滑丛同态，则 $F$ 为光滑丛同构。
\end{proposition}

\begin{exercise}[10.27]
证明：秩为 $k$ 的光滑向量丛光滑平凡当且仅当它在 $M$ 上光滑同构于积丛 $M\times\mathbb{R}^k$。
\end{exercise}

\begin{example}[10.28（丛同态的例子）]
\begin{enumerate}[label=(\alph*)]
\item 若 $F:M\to N$ 光滑，则微分 $dF:TM\to TN$ 是覆盖 $F$ 的光滑丛同态。
\item 若 $E\to M$ 光滑向量丛且 $S\subseteq M$ 为浸入子流形，则包含 $E|_S\hookrightarrow E$ 是覆盖 $S\hookrightarrow M$ 的光滑丛同态。
\end{enumerate}
\end{example}

\subsection{截面层面的刻画（$C^\infty(M)$-线性）}

\begin{definition}[$C^\infty(M)$-线性]
设 $E,E'\to M$ 为光滑向量丛。映射 $\mathcal{F}:\Gamma(E)\to \Gamma(E')$ 称\textbf{关于 $C^\infty(M)$ 线性}，若对任意 $u_1,u_2\in C^\infty(M)$ 与 $\sigma_1,\sigma_2\in\Gamma(E)$，
\[
\mathcal{F}(u_1\sigma_1+u_2\sigma_2)=u_1\mathcal{F}(\sigma_1)+u_2\mathcal{F}(\sigma_2).
\]
\end{definition}

\begin{lemma}[10.29（丛同态刻画引理）]
设 $E,E'\to M$ 为光滑向量丛。映射 $\mathcal{F}:\Gamma(E)\to \Gamma(E')$ 关于 $C^\infty(M)$ 线性，当且仅当存在在 $M$ 上的光滑丛同态 $F:E\to E'$ 使得对所有 $\sigma\in\Gamma(E)$，
\[
\mathcal{F}(\sigma)=F\circ \sigma.
\]
\end{lemma}

\begin{example}[10.30（流形上的丛同态：由 $C^\infty$-线性给出）]
\begin{enumerate}[label=(\alph*)]
\item 若 $f\in C^\infty(M)$，映射 $\mathfrak{X}(M)\to\mathfrak{X}(M)$，$X\mapsto fX$ 关于 $C^\infty(M)$ 线性，因此给出 $TM\to TM$ 的丛同态。
\item 若 $Z$ 为 $\mathbb{R}^3$ 上光滑向量场，则 $X\mapsto X\times Z$ 关于 $C^\infty(\mathbb{R}^3)$ 线性，给出 $T\mathbb{R}^3\to T\mathbb{R}^3$ 的丛同态。
\item 若 $Z\in\mathfrak{X}(\mathbb{R}^n)$，点积 $X\mapsto X\cdot Z$ 关于 $C^\infty(\mathbb{R}^n)$ 线性，给出 $T\mathbb{R}^n\to \mathbb{R}^n\times\mathbb{R}$ 的丛同态。
\end{enumerate}
\end{example}

\section{子丛（Subbundles）}

\begin{definition}[子丛与光滑子丛]
设 $\pi_E:E\to M$ 为向量丛。若 $D\subseteq E$ 为子空间，$\pi_D=\pi_E|_D:D\to M$ 使得：
\begin{itemize}
\item 对每个 $p\in M$，$D_p=D\cap E_p$ 是 $E_p$ 的线性子空间；
\item $D\to M$ 本身为向量丛（因此各 $D_p$ 非空且维数相同），
\end{itemize}
则称 $D\to M$ 为 $E$ 的\textbf{子丛}。
若 $E$ 光滑且 $D$ 既是光滑向量丛又作为子集是 $E$ 的嵌入子流形，则称为\textbf{光滑子丛}。
\end{definition}

\begin{exercise}[10.31]
设 $E\to M$ 光滑向量丛，$D\subseteq E$ 为光滑子丛。证明包含 $\iota:D\hookrightarrow E$ 是在 $M$ 上的光滑丛同态。
\end{exercise}

\begin{lemma}[10.32（子丛的局部标架判别）]
设 $E\to M$ 为光滑向量丛，并对每个 $p\in M$ 给定 $m$ 维线性子空间 $D_p\subseteq E_p$，令
\[
D=\bigcup_{p\in M}D_p\subseteq E.
\]
则 $D$ 是 $E$ 的光滑子丛当且仅当：对每个 $p\in M$，存在邻域 $U$ 以及光滑局部截面 $\sigma_1,\dots,\sigma_m:U\to E$，使得对每个 $q\in U$，
\[
\{\sigma_1(q),\dots,\sigma_m(q)\}\ \text{是 }D_q\text{ 的一组基。}
\]
\end{lemma}

\begin{example}[10.33（子丛的例子）]
\begin{enumerate}[label=(\alph*)]
\item 若 $V$ 为 $M$ 上处处非零的光滑向量场，则由每点的 $\mathrm{span}(V_p)\subset T_pM$ 得到 $TM$ 的秩 $1$ 光滑子丛。
\item 若 $E\to M$ 平凡且有光滑整体标架 $(E_1,\dots,E_k)$，令 $D_p=\mathrm{span}(E_1|_p,\dots,E_m|_p)$，则 $D$ 为光滑子丛。
\item 若 $S\subseteq M$ 为浸入 $k$ 维子流形，则 $TS$ 是环境切丛 $TM|_S$ 的秩 $k$ 光滑子丛（与练习 10.14 相关）。
\end{enumerate}
\end{example}

\subsection{常秩丛同态的核与像}

\begin{theorem}[10.34（常秩 $\Leftrightarrow$ 核与像为子丛）]
设 $E,E'\to M$ 为光滑向量丛，$F:E\to E'$ 为在 $M$ 上的光滑丛同态。定义
\[
\ker F=\bigcup_{p\in M}\ker(F|_{E_p})\subseteq E,\qquad
\mathrm{Im}\,F=\bigcup_{p\in M}\mathrm{Im}(F|_{E_p})\subseteq E'.
\]
则 $\ker F$ 与 $\mathrm{Im}\,F$ 分别是 $E$ 与 $E'$ 的光滑子丛，当且仅当 $F$ 具有\textbf{常秩}（即 $\mathrm{rank}(F|_{E_p})$ 与 $p$ 无关）。
\end{theorem}

\begin{lemma}[10.35（正交补丛）]
设 $M\subseteq \mathbb{R}^n$ 为浸入子流形，$D$ 为 $T\mathbb{R}^n|_M$ 的秩 $k$ 光滑子丛。
对每个 $p\in M$，令 $D_p^\perp\subseteq T_p\mathbb{R}^n$ 为欧氏内积下的正交补，并定义
\[
D^\perp=\{(p,v)\in T\mathbb{R}^n:\ p\in M,\ v\in D_p^\perp\}\subseteq T\mathbb{R}^n|_M.
\]
则 $D^\perp$ 是秩 $n-k$ 的光滑子丛；且每点附近存在 $D^\perp$ 的光滑正交归一局部标架。
\end{lemma}

\begin{corollary}[10.36（$\mathbb{R}^n$ 中子流形的法丛）]
若 $M\subseteq \mathbb{R}^n$ 为浸入的 $m$ 维子流形，则其\textbf{法丛} $NM$ 是 $T\mathbb{R}^n|_M$ 的秩 $n-m$ 光滑子丛，且每点附近存在 $NM$ 的光滑正交归一局部标架。
\end{corollary}

\section{纤维丛（Fiber Bundles）：向量丛的推广}

\begin{definition}[纤维丛（Fiber bundle）]
设 $M,F$ 为拓扑空间。一个\textbf{以 $F$ 为模型纤维}、定义在 $M$ 上的\textbf{纤维丛}是拓扑空间 $E$ 连同满射连续映射
\[
\pi:E\to M
\]
使得对每个 $x\in M$，存在邻域 $U\ni x$ 与同胚
\[
\Phi:\pi^{-1}(U)\to U\times F
\]
（局部平凡化）满足 $\pi_1\circ\Phi=\pi$。
若 $E,M,F$ 为光滑流形、$\pi$ 光滑且可取 $\Phi$ 为微分同胚，则称为\textbf{光滑纤维丛}。
若存在全局平凡化（整个 $M$ 上的一个局部平凡化），则称为\textbf{平凡纤维丛}（光滑情形为\textbf{光滑平凡}）。
\end{definition}

\begin{example}[10.37（纤维丛的例子）]
\begin{enumerate}[label=(\alph*)]
\item 任意积空间 $M\times F$ 连同投影 $\pi_1$ 是纤维丛（积纤维丛），且平凡。
\item 任意秩 $k$ 向量丛都是以 $\mathbb{R}^k$ 为模型纤维的纤维丛。
\item 莫比乌斯丛中 $\mathbb{R}\times[-1,1]$ 在商映射下的像给出一个底空间为 $S^1$、纤维为 $[-1,1]$ 的纤维丛（莫比乌斯带），它不是平凡纤维丛。
\item 任意覆盖映射 $\pi:E\to M$ 都是模型纤维为离散空间的纤维丛。
\end{enumerate}
\end{example}

\end{document}